% ============================================================
% SECCIÓN 5: Backend Centralizado
% ============================================================
\chapter{Servidor Central: Backend en Google Cloud}
\label{sec:backend}

% ─────────────────────────────────────────────────────────────
\section{Visión General del Servidor}

El \textbf{servidor central} de Demeter es el núcleo digital del sistema: el punto de convergencia
donde la telemetría de campo adquiere persistencia, donde los experimentos científicos cobran forma, 
y desde donde los investigadores pueden actuar sobre el entorno físico del invernadero. Está 
desplegado en una \textbf{máquina virtual de Google Cloud Platform} (GCP), con acceso público 
protegido por TLS y orientado a múltiples clientes simultáneos.

El software del servidor está organizado en un conjunto de \textbf{contenedores Docker independientes} 
que colaboran a través de una red interna privada, sin puertos expuestos entre ellos al exterior, 
siguiendo el principio de mínimo acceso. Un \textbf{Nginx} actúa como proxy inverso único, 
terminando TLS y dirigiendo el tráfico al servicio correcto.

La arquitectura completa de contenedores, volúmenes y comunicaciones se muestra en la 
Sección~\ref{fig:anx_backend_docker} del Anexo.

% ─────────────────────────────────────────────────────────────
\section{Arquitectura de Contenedores Docker}

La pila de producción del servidor se compone de \textbf{seis contenedores} que conviven en una 
red privada interna denominada \texttt{demeter-net}:

\begin{table}[H]
    \centering
    \caption{Contenedores del sistema Demeter en producción}
    \begin{tabular}{lllp{5.8cm}}
        \toprule
        \textbf{Contenedor} & \textbf{Imagen base} & \textbf{Escala} & \textbf{Función} \\
        \midrule
        \texttt{demeter-frontend}   & Nginx + React build  & 1 instancia   & SPA distribuida al navegador. \\
        \texttt{demeter-api}        & Python (FastAPI)     & 1 instancia   & API REST + WebSocket + Lógica. \\
        \texttt{demeter-math-worker}& Python (misma base)  & Escalable     & Worker de cálculos pesados. \\
        \texttt{demeter-sequencer}  & Python (misma base)  & 1 instancia   & Motor de secuencias temporizadas. \\
        \texttt{demeter-db}         & PostgreSQL 15        & 1 instancia   & Persistencia relacional. \\
        \texttt{demeter-redis}      & Redis Alpine         & 1 instancia   & Broker, caché y cola de tareas. \\
        \bottomrule
    \end{tabular}
    \label{tab:docker_containers}
\end{table}

\textbf{Volúmenes persistentes:}
\begin{itemize}[leftmargin=*, label=--]
    \item \texttt{postgres\_data}: datos de la base de datos PostgreSQL. 
    \item \texttt{shared\_data}: archivos intermedios generados por el worker 
          y accesibles desde la API (\texttt{/app/shared}).
    \item \texttt{export\_data}: archivos de exportación de datos históricos 
          (\texttt{/app/exports}).
\end{itemize}

\textbf{Variables de entorno}: cada contenedor recibe su configuración mediante un fichero 
\texttt{.env}, evitando credenciales en el repositorio de código. La URL de Redis, PostgreSQL, 
los secretos JWT y las claves de API son inyectadas en tiempo de ejecución.


% ─────────────────────────────────────────────────────────────
\section{Redis: Bus de Mensajes y Caché}

Redis desempeña el rol más transversal del sistema: no es únicamente una caché de velocidad, 
sino el \textbf{bus de mensajes central} que cohesiona los distintos procesos Python que corren 
en contenedores independientes. Opera con dos bases de datos lógicas:

\begin{table}[H]
    \centering
    \caption{Bases de datos Redis y su propósito}
    \begin{tabular}{llp{8cm}}
        \toprule
        \textbf{DB} & \textbf{URL} & \textbf{Datos} \\
        \midrule
        DB/0 & \texttt{redis://...6379/0} & Caché de aplicación y Pub/Sub \\
        DB/1 & \texttt{redis://...6379/1} & Cola de tareas RQ (\textit{demeter\_tasks}) \\
        \bottomrule
    \end{tabular}
\end{table}

\subsection{Las Cuatro Categorías de Datos en Redis}

\begin{enumerate}

\item \textbf{Estado de Dispositivo} (\texttt{device:\{id\}:state})\\ 
Almacena el estado ON/OFF \textit{confirmado} de cada pin de actuador, tras recibir la 
confirmación de ejecución del nodo (CMD\_ACK). Persiste indefinidamente hasta ser sobreescrito.

\item \textbf{Caché de Intención} (\texttt{demeter:cache:device:\{id\}}, TTL 60~s)\\
Cuando el usuario pulsa un botón, el estado \textit{previsto} se almacena inmediatamente con 
un TTL de 60 segundos. El dashboard refleja este estado optimista antes de recibir confirmación 
del hardware, haciendo la interfaz percibida como instantánea.

\item \textbf{Telemetría} (\texttt{sensor:\{id\}:telemetry}, TTL 60~s)\\
Cada nueva trama de sensor que llega desde la Raspberry Pi es almacenada en Redis con expiración 
de un minuto. Endpoints de tiempo real la consultan directamente para responder en 
microsegundos, sin tocar PostgreSQL.

\item \textbf{Lote de Actividad} (\texttt{demeter:activity:batch}, cola FIFO)\\
Cada acción del usuario (botón pulsado, secuencia ejecutada) se encola en Redis en lugar de 
escribirse directamente en PostgreSQL. Una tarea en background (\textit{Batch Flusher}) 
los vuelca periódicamente en la tabla \texttt{ActivityLog} de forma agrupada, evitando una 
avalancha de escrituras en la base de datos ante interacciones rápidas.

\end{enumerate}

\subsection{Canales Pub/Sub}

Redis Pub/Sub es el mecanismo que permite a procesos en diferentes contenedores comunicarse 
sin acoplamiento directo:

\begin{table}[H]
    \centering
    \caption{Canales Pub/Sub y sus productores/consumidores}
    \begin{tabular}{lll}
        \toprule
        \textbf{Canal} & \textbf{Publicador} & \textbf{Suscriptor} \\
        \midrule
        \texttt{demeter:commands} & \texttt{POST /api/command} & \texttt{WS Dispatcher} (asyncio task) \\
        \texttt{iot\_events}      & RedisManager               & Frontend (retransmisión WS broadcast) \\
        \bottomrule
    \end{tabular}
\end{table}

El flujo completo de un comando desde el navegador hasta el ESP32, pasando por Redis y la 
Raspberry Pi, se muestra en la Sección~\ref{fig:anx_redis_flow} del Anexo.


% ─────────────────────────────────────────────────────────────
\section{PostgreSQL: Persistencia y Esquema de Datos}

PostgreSQL almacena toda la información de largo plazo del sistema. El esquema evoluciona 
mediante \textbf{Alembic}, el gestor de migraciones de SQLAlchemy, garantizando que cualquier 
cambio en el modelo de datos sea reproducible y trazable en todos los entornos.

Las tablas principales son:

\begin{table}[H]
    \centering
    \caption{Tablas principales de PostgreSQL}
    \begin{tabular}{lp{10cm}}
        \toprule
        \textbf{Tabla}            & \textbf{Datos almacenados} \\
        \midrule
        \texttt{users}            & Credenciales de acceso, roles (admin/viewer) y estado activo. \\
        \texttt{plants}           & Registro botánico de plantas: especie, variedad, fecha de siembra, estado vital, metadatos científicos, nodo asignado. \\
        \texttt{experiments}      & Agrupación de plantas en experimentos de investigación. \\
        \texttt{experimento\_planta\_link} & Tabla intermedia N:M entre plantas y experimentos. \\
        \texttt{telemetry\_th}    & Serie temporal de mediciones de temperatura y humedad por nodo. \\
        \texttt{activity\_log}    & Registro de acciones (comandos enviados, interacciones de usuario). \\
        \texttt{sequences}        & Secuencias de pasos de riego configuradas por los usuarios. \\
        \bottomrule
    \end{tabular}
    \label{tab:postgres_tables}
\end{table}

El diagrama entidad-relación completo puede consultarse en la Sección~\ref{fig:anx_db_schema} del Anexo.


% ─────────────────────────────────────────────────────────────
\section{El Worker Matemático}

El \textbf{Math Worker} (\texttt{demeter-math-worker}) es un proceso Python independiente que 
escucha la cola de tareas \texttt{demeter\_tasks} en Redis (DB/1) usando la librería 
\textbf{RQ} (\textit{Redis Queue}).

Su función es absorber los cálculos \textit{pesados y no urgentes} que no deben bloquear la 
API principal. Las tareas que ejecuta incluyen:

\begin{itemize}[leftmargin=*, label=--]
    \item \textbf{Análisis científico por experimento}: cálculo de estadísticos de las series 
          temporales de temperatura y humedad asociadas a un experimento (media, varianza, 
          tendencias, correlación entre nodos).
    \item \textbf{Generación de informes}: empaquetamiento de los análisis en archivos de 
          resultados exportables al volumen compartido \texttt{/app/shared}.
    \item \textbf{Nightly Batch}: tareas programadas de limpieza, compactación de telemetría 
          histórica y recálculo de estadísticos agregados.
\end{itemize}

El proceso de encolar una tarea sigue este patrón: el cliente llama a 
\texttt{POST /api/analysis/generate}, el backend encola la tarea en Redis y devuelve un 
\texttt{task\_id}. El cliente puede consultar \texttt{GET /api/analysis/status/\{task\_id\}} 
para conocer el progreso y, una vez completada, descarga el resultado de 
\texttt{GET /api/analysis/download/\{filename\}}.


% ─────────────────────────────────────────────────────────────
\section{El Secuenciador}

El \textbf{Secuenciador} (\texttt{demeter-sequencer}) es un proceso Python independiente con 
un único propósito: ejecutar secuencias de pasos de riego multi-acción de forma temporizada, 
sin bloquear la API ni consumir recursos del worker matemático.

Una secuencia puede definir, por ejemplo: \textit{``Abrir válvula A durante 10 segundos, 
esperar 2 segundos, abrir bomba B durante 30 segundos, cerrar todo''}. 

El Secuenciador traduce cada paso temporizado en un comando individual y lo lanza al canal 
Redis \texttt{demeter:commands}, que el Dispatcher convierte en una trama que llega al 
actuador físico. Entre paso y paso, duerme el tiempo de espera definido sin bloquear 
ningún otro proceso del sistema.


% ─────────────────────────────────────────────────────────────
\section{El Backend FastAPI}

La \textbf{API principal} (\texttt{demeter-api}) es el componente central y más complejo 
del servidor. Implementada con el framework \textbf{FastAPI} sobre \texttt{asyncio}, 
combina en un único proceso el servidor HTTP REST, el servidor WebSocket, 
y dos tareas en background permanentes:

\begin{itemize}[leftmargin=*, label=--]
    \item \textbf{Redis Command Listener}: escucha el canal \texttt{demeter:commands} y 
          retransmite comandos a la Raspberry Pi conectada por WebSocket.
    \item \textbf{Activity Batch Flusher}: cada 30 segundos vuelca el lote de 
          eventos de actividad de Redis a PostgreSQL.
\end{itemize}

Todas las rutas que acceden a la base de datos lo hacen de forma \textbf{completamente asíncrona} 
mediante \texttt{AsyncSession} de SQLAlchemy, garantizando que las consultas lentas no 
bloqueen otras peticiones en curso.


% ─────────────────────────────────────────────────────────────
\section{El Frontend React}

El \textbf{frontend} es una \textit{Single Page Application} (SPA) construida en React,
servida como ficheros estáticos a través de Nginx. Una vez cargada en el navegador, 
se comunica directamente con la API del backend mediante peticiones HTTPS y WebSocket. 
Ofrece las siguientes funcionalidades:

\begin{table}[H]
    \centering
    \caption{Módulos del frontend React}
    \begin{tabular}{lp{9.5cm}}
        \toprule
        \textbf{Módulo}          & \textbf{Funcionalidades} \\
        \midrule
        Dashboard en tiempo real & Visualización de temperatura, humedad y suelo con gráficos actualizados vía WebSocket. \\
        Control Manual           & Botones para activar/desactivar pines de actuadores individualmente. \\
        Planificador de Secuencias & Creación, guardado y ejecución de secuencias de riego multi-paso. \\
        LIMS                     & Registro y gestión de plantas, asignación a experimentos, consulta de telemetría por planta. \\
        SDK / Exportación        & Descarga de datos en formatos académicos, acceso programático para análisis externos. \\
        \bottomrule
    \end{tabular}
\end{table}


% ─────────────────────────────────────────────────────────────
\section{Documentación de la API REST}

El servidor expone una API REST bajo el prefijo \texttt{/api}, con documentación interactiva 
generada automáticamente por FastAPI en \texttt{/api/docs}. A continuación se detallan todos 
los \textit{routers} y sus endpoints más relevantes.

\subsection{Sistema (\texttt{/api/system})}
\begin{table}[H]
    \centering
    \begin{tabular}{llp{7.5cm}}
        \toprule
        \textbf{Método} & \textbf{Ruta}   & \textbf{Descripción} \\
        \midrule
        \texttt{GET}  & \texttt{/api/system/health} & Comprueba el estado del servidor y la conexión a Redis y PostgreSQL. Devuelve un JSON de estado. \\
        \bottomrule
    \end{tabular}
\end{table}

\subsection{Autenticación (\texttt{/api/auth})}
\begin{table}[H]
    \centering
    \begin{tabular}{llp{7.5cm}}
        \toprule
        \textbf{Método} & \textbf{Ruta}             & \textbf{Descripción} \\
        \midrule
        \texttt{POST} & \texttt{/api/auth/login}    & Recibe usuario y contraseña. Devuelve un JWT Bearer Token con el rol del usuario y expiración configurable. \\
        \texttt{POST} & \texttt{/api/auth/register} & Registra un nuevo usuario (admin). Almacena la contraseña hasheada con bcrypt. \\
        \texttt{GET}  & \texttt{/api/auth/me}       & Devuelve el perfil del usuario activo autenticado mediante el token. \\
        \bottomrule
    \end{tabular}
\end{table}

\subsection{Comandos (\texttt{/api/command})}
\begin{table}[H]
    \centering
    \begin{tabular}{llp{7.5cm}}
        \toprule
        \textbf{Método} & \textbf{Ruta}          & \textbf{Descripción} \\
        \midrule
        \texttt{POST} & \texttt{/api/command}    & Envía un comando al hardware. El campo \texttt{type} discrimina el subtipo (\texttt{set\_gpio}, \texttt{set\_pwm}, \texttt{exec\_sequence}, \texttt{ping}, \texttt{get\_sensors}). El comando es validado con Pydantic y publicado en Redis. Devuelve \texttt{202 Accepted}. \\
        \bottomrule
    \end{tabular}
\end{table}

\subsection{Descubrimiento (\texttt{/api/discovery})}
\begin{table}[H]
    \centering
    \begin{tabular}{llp{7.5cm}}
        \toprule
        \textbf{Método} & \textbf{Ruta}               & \textbf{Descripción} \\
        \midrule
        \texttt{GET}  & \texttt{/api/discovery/nodes} & Lista todos los nodos activos conocidos por el sistema con su último estado. \\
        \texttt{GET}  & \texttt{/api/discovery/nodes/\{id\}} & Estado detallado y telemetría en caché de un nodo específico. \\
        \bottomrule
    \end{tabular}
\end{table}

\subsection{Historial de Telemetría (\texttt{/api/history})}
\begin{table}[H]
    \centering
    \begin{tabular}{llp{7.5cm}}
        \toprule
        \textbf{Método} & \textbf{Ruta}              & \textbf{Descripción} \\
        \midrule
        \texttt{GET}  & \texttt{/api/history/node/\{id\}} & Devuelve hasta \texttt{days} días de telemetría histórica de temperatura y humedad para el nodo indicado. Requiere autenticación JWT. \\
        \bottomrule
    \end{tabular}
\end{table}

\subsection{LIMS --- Plantas y Experimentos (\texttt{/api/lims})}
\begin{table}[H]
    \centering
    \begin{tabular}{llp{6.8cm}}
        \toprule
        \textbf{Método}   & \textbf{Ruta}                   & \textbf{Descripción} \\
        \midrule
        \texttt{GET}    & \texttt{/api/lims/plantas}          & Lista todas las plantas registradas. \\
        \texttt{POST}   & \texttt{/api/lims/plantas}          & Registra una nueva planta con metadatos botánicos completos. \\
        \texttt{GET}    & \texttt{/api/lims/plantas/\{id\}}   & Detalle de una planta con sus experimentos asociados. \\
        \texttt{PATCH}  & \texttt{/api/lims/plantas/\{id\}}   & Actualiza parcialmente los datos de una planta. \\
        \texttt{GET}    & \texttt{/api/lims/plantas/\{id\}/telemetry} & Telemetría de 6 meses del nodo asignado a la planta, directo de Redis. \\
        \texttt{GET}    & \texttt{/api/lims/experimentos}     & Lista todos los experimentos vigentes. \\
        \texttt{POST}   & \texttt{/api/lims/experimentos/generar} & Crea un experimento y asigna plantas existentes. \\
        \bottomrule
    \end{tabular}
\end{table}

\subsection{Análisis Científico (\texttt{/api/analysis})}
\begin{table}[H]
    \centering
    \begin{tabular}{llp{6.8cm}}
        \toprule
        \textbf{Método} & \textbf{Ruta}                     & \textbf{Descripción} \\
        \midrule
        \texttt{POST} & \texttt{/api/analysis/generate}       & Encola un análisis estadístico completo de un experimento en el Math Worker (RQ). Devuelve el \texttt{task\_id}. \\
        \texttt{GET}  & \texttt{/api/analysis/status/\{id\}}  & Consulta el estado de la tarea encolada: \textit{queued}, \textit{started}, \textit{finished} o \textit{failed}. \\
        \texttt{GET}  & \texttt{/api/analysis/download/\{f\}} & Descarga el archivo de análisis generado (ZIP) desde el volumen compartido. \\
        \bottomrule
    \end{tabular}
\end{table}

\subsection{SDK Científico (\texttt{/api/sdk})}
\begin{table}[H]
    \centering
    \begin{tabular}{llp{7.5cm}}
        \toprule
        \textbf{Método} & \textbf{Ruta}                     & \textbf{Descripción} \\
        \midrule
        \texttt{GET}  & \texttt{/api/sdk/mediciones/\{exp\_id\}} & Devuelve la telemetría cruda en JSON de todas las plantas del experimento. Protegido por \texttt{X-API-Key}. Prioriza caché Redis sobre consulta directa a PostgreSQL. \\
        \bottomrule
    \end{tabular}
\end{table}

\subsection{Exportación (\texttt{/api/export})}
\begin{table}[H]
    \centering
    \begin{tabular}{llp{7.5cm}}
        \toprule
        \textbf{Método} & \textbf{Ruta}             & \textbf{Descripción} \\
        \midrule
        \texttt{GET}  & \texttt{/api/export/\{formato\}} & Exporta los datos históricos en formatos CSV o XLSX para análisis externo. Se genera y devuelve como descarga directa. \\
        \bottomrule
    \end{tabular}
\end{table}


% ─────────────────────────────────────────────────────────────
\section{El Canal WebSocket}

El servidor expone el endpoint \texttt{/ws/\{client\_id\}} como punto de entrada único para 
todas las conexiones WebSocket. El comportamiento varía según el identificador con el que 
se conecta el cliente:

\begin{table}[H]
    \centering
    \caption{Diferencia de comportamiento por tipo de cliente WebSocket}
    \begin{tabular}{lllp{5cm}}
        \toprule
        \textbf{client\_id}          & \textbf{Tipo}    & \textbf{Dirección} & \textbf{Mensajes} \\
        \midrule
        \texttt{raspberry\_gateway}  & Máquina          & Bidireccional      & Recibe JSON de comandos, envía JSON de telemetría. \\
        Cualquier otro               & Frontend/Monitor & Solo recepción     & Solo recibe push de telemetría. Los intentos de enviar comandos son rechazados con un mensaje informativo. \\
        \bottomrule
    \end{tabular}
\end{table}

Esta separación es una decisión de diseño deliberada: el canal WebSocket queda reservado para 
el flujo de datos en tiempo real (telemetría), mientras los comandos circulan por HTTP POST. 
Esto simplifica el frontend (sin lógica de reconexión de socket) y hace los comandos 
trivialmente verificables con herramientas estándar como \texttt{curl} o la documentación 
Swagger de la API.
